\chapter{Очакване, дисперсия, пораждащи функции}

\section{Математическо очакване}

\begin{defbox}[~(Математическо очакване)]
Нека $X$ е дискретна случайна величина. Ако редът $\sum_i x_i p_i$ е
добре дефиниран (т.е.\ крайна сума или абсолютно сходящ ред), то
\[
  \mathbb E X = \sum_i x_i p_i = \sum_i x_i\,\mathbb P(X=x_i)
\]
се нарича \textbf{математическо очакване} на $X$. Ако $X$ приема
краен брой стойности, $\mathbb E X$ винаги съществува.
\end{defbox}

\begin{rembox}[Частен случай: равномерно разпределение]
Ако $X$ е равномерно разпределена върху крайното множество
$\{x_1,\dots,x_n\}$ (т.е.\ $\mathbb P(X=x_i)=\tfrac1n$), то
\[
  \mathbb E X = \sum_{i=1}^{n} x_i\cdot\frac1n = \frac{\sum_{i=1}^{n}x_i}{n},
\]
което е точно \emph{средноаритметичното} на стойностите, които $X$
може да приема.
\end{rembox}

\begin{rembox}[Забележка: очакването не винаги съществува]
Възможно е редът $\sum_i |x_i|\,p_i$ да се разлива дори когато
$\sum_i p_i=1$ коректно определя разпределение. Класически пример е
разпределението $\mathbb P(X=j)=\tfrac{6}{\pi^2}\cdot\tfrac1{j^2}$ за
$j\ge 1$ (коректно, защото $\sum \tfrac1{j^2}=\tfrac{\pi^2}{6}$), за
което $\mathbb E X = \tfrac{6}{\pi^2}\sum_{j\ge1}\tfrac1j = \infty$ -
редът на средните стойности се разлива, следователно $X$ няма
математическо очакване.
\end{rembox}

\section{Свойства на математическото очакване}

\begin{thmbox}[]{Лема}
Нека $X$ и~$Y$ са две дискретни случайни величини, за които
$\mathbb E X$ и~$\mathbb E Y$ съществуват ($\mathbb EX<\infty$,
$\mathbb EY<\infty$). Тогава:
\begin{enumerate}[label=\arabic*)]
  \item \textbf{Очакване на константа.} Ако $X\equiv c$, то $\mathbb E X = c$.
  \item \textbf{Очакване на индикаторна случайна величина.} Ако $X=I_A$, то $\mathbb EX = \mathbb P(A)$ за някое събитие $A$.
  \item \textbf{Умножение с~константа.} За $c\in\mathbb R$: $\mathbb E(cX) = c\,\mathbb EX$.
  \item \textbf{Линейност.} $\mathbb E(aX+bY) = a\,\mathbb EX + b\,\mathbb EY$ (важи и~без константите); по-общо, за краен брой случайни величини: $\mathbb E\sum_{i=1}^n c_iX_i = \sum_{i=1}^n c_i\,\mathbb EX_i$.
  \item \textbf{Монотонност.} Ако $X\le Y$ (с~вероятност~1), то $\mathbb EX\le\mathbb EY$.
  \item \textbf{Очакване на функция на случайна величина.} Ако $g:\mathbb R\to\mathbb R$ е измерима функция, то за дискретна $X$: $\mathbb Eg(X) = \sum_{x\in\mathcal X} g(x)\,p_X(x)$, при условие че редът е абсолютно сходящ.
  \item \textbf{Независимост.} Ако $X,Y$ са независими случайни величини, то $\mathbb E(XY)=\mathbb EX\cdot\mathbb EY$ (обратното не е вярно в~общия случай).
  \item \textbf{Неравенство на Йенсен.} Ако $\varphi:\mathbb R\to\mathbb R$ е изпъкнала функция и~$X$ е случайна величина с~крайно очакване, то $\varphi(\mathbb EX)\le\mathbb E(\varphi(X))$.
\end{enumerate}
\end{thmbox}

\begin{proofbox}
\textbf{1)} $\mathbb EX = c\cdot\mathbb P(c=c)=c\cdot1=c$.

\textbf{2)} Нека $A\in\mathcal A$ е събитие и~$I_A(\omega)=1$
за $\omega\in A$, $0$ за $\omega\notin A$. Тогава $\mathbb E\mathbb
1_A = 1\cdot\mathbb P(A) + 0\cdot\mathbb P(A^c) = \mathbb P(A)$.

\textbf{3)} $Z=cX=g(X)$ за $g(x)=cx$, следователно $\mathbb EZ =
\mathbb E(cX) = \sum_j g(x_j)p_j = \sum_j cx_j p_j = c\sum_j x_j p_j =
c\,\mathbb EX$.

\textbf{4)} Достатъчно е да покажем $\mathbb E(X+Y)=\mathbb EX+
\mathbb EY$ (общият случай с~константи следва от~3). Нека $p_{ij}=
\mathbb P(X=x_i,Y=y_j)$. Тогава
\begin{align*}
  \mathbb E(X+Y) &= \sum_{i,j}(x_i+y_j)p_{ij} = \sum_i x_i\sum_j p_{ij} + \sum_j y_j\sum_i p_{ij}\\
  &= \sum_i x_i\,\mathbb P(X=x_i) + \sum_j y_j\,\mathbb P(Y=y_j) = \mathbb EX+\mathbb EY,
\end{align*}
защото $\mathbb P(X=x_i)=\sum_j\mathbb P(X=x_i,Y=y_j)=\sum_j p_{ij}$
- сумирайки съвместните вероятности по всички стойности на $Y$, ние
покриваме всички сценарии за $Y$ и~получаваме пълната (маргинална)
вероятност $X$ да приеме стойността $x_i$, независимо каква стойност
е приело~$Y$.

\textbf{5)} Нека $X\le Y$ с~вероятност~1, т.е.\ за всяка двойка
$(x,y)$ с~$x>y$ имаме $p_{X,Y}(x,y)=0$. За простота, ако $X,Y$ са
дефинирани върху едно и~също вероятностно пространство, разглеждаме
$Z=Y-X\ge0$ с~вероятност~1. Тогава $\mathbb EZ = \sum_{z\ge0} z\,
p_Z(z)\ge 0$, защото всички $z\ge0$ и~вероятностите са неотрицателни.
От линейността (4): $\mathbb EY-\mathbb EX=\mathbb E(Y-X)\ge 0
\Rightarrow \mathbb EX\le\mathbb EY$.

\textbf{6)} Нека $X$ приема стойности в~изброимото множество
$\mathcal X$ с~$p_X(x)=\mathbb P(X=x)$, т.е.\ $X=\sum_{x_i\in
\mathcal X} x_i\,I_{H_i}$, където $H_i=\{X=x_i\}$. От
теоремата за смяна на променливите (Тема~4), $g(X) = \sum_i
g(x_i)\,I_{H_i}$ е дискретна случайна величина, приемаща
стойност $g(x_i)$ върху $H_i$ (при съвпадащи стойности на $g$ за
различни $x_i$, съответните $H_i$ просто се обединяват). Прилагайки
дефиницията на очакване към тази нова дискретна случайна величина:
\[
  \mathbb Eg(X) = \sum_i g(x_i)\,\mathbb P(H_i) = \sum_{x\in\mathcal X} g(x)\,p_X(x),
\]
при условие че този ред е абсолютно сходящ (нужно е сумата да не
зависи от реда на събиране, тъй като при групиране на равни стойности
$g(x_i)$ редът на сумиране естествено се променя). \qedbg

\textbf{7)} Нека $g(x,y)=xy$. Тогава, ползвайки независимостта
$X\perp\!\!\!\perp Y$:
\begin{align*}
  \mathbb E(XY) &= \sum_{i,j} x_iy_j p_{ij} \overset{X\perp\!\!\!\perp Y}{=} \sum_{i,j} x_iy_j\,\mathbb P(X=x_i)\,\mathbb P(Y=y_j)\\
  &= \sum_i x_i\,\mathbb P(X=x_i)\sum_j y_j\,\mathbb P(Y=y_j) = \mathbb EX\cdot\mathbb EY.
\end{align*}

\textbf{8) (Неравенство на Йенсен)} Нека $\mu=\mathbb EX$. Понеже
$\varphi$ е изпъкнала, в~точката $\mu$ съществува \emph{опорна права}
$\ell(x)=\varphi(\mu)+c(x-\mu)$ (с~$c$ равно на производната
$\varphi'(\mu)$, ако тя съществува, или произволен subgradient в~противен случай), такава че $\varphi(x)\ge \ell(x)$ за всяко
$x\in\mathbb R$ (това е характеристично свойство на изпъкналите
функции - графиката им лежи над всяка допирателна права). Прилагаме
монотонността (5) и~линейността (4) към неравенството
$\varphi(X)\ge\ell(X)=\varphi(\mu)+c(X-\mu)$:
\[
  \mathbb E(\varphi(X)) \ge \mathbb E(\ell(X)) = \varphi(\mu) + c\big(\mathbb EX-\mu\big) = \varphi(\mu) + c\cdot 0 = \varphi(\mathbb EX). \qedbg
\]
\end{proofbox}

\section{Дисперсия}

\begin{rembox}[Идея]
Дисперсията е мярка за разпространението (разсейването) на
стойностите на случайна величина спрямо нейната средна стойност
(математическото очакване). Тя показва количествено колко стойностите
се отклоняват от средната.
\end{rembox}

\begin{defbox}[~(Дисперсия)]
Нека $X$ е дискретна случайна величина и~нека сумата
$\sum_j (x_j-\mathbb EX)^2 p_j$ е добре дефинирана (т.е.\ крайна).
Тогава
\[
  DX = \mathbb E(X-\mathbb EX)^2 = \sum_j (x_j-\mathbb EX)^2 p_j
\]
се нарича \textbf{дисперсия} на $X$. Това е средноаритметичното на
квадратите на отклоненията на всяка стойност от средната.
\end{defbox}

\begin{defbox}[~(Стандартно отклонение)]
Ако $X$ е случайна величина и~$DX$ е крайно, то $\sqrt{DX}$ се нарича
\textbf{стандартно отклонение} на $X$.
\end{defbox}

\begin{thmbox}[]{Теорема}
Нека $X$ е дискретна случайна величина и~нека $DX<\infty$. Тогава
\[
  DX = \mathbb E(X-\mathbb EX)^2 = \mathbb EX^2 - (\mathbb EX)^2.
\]
\end{thmbox}

\begin{proofbox}
\[
  \mathbb E(X-\mathbb EX)^2 = \mathbb E\big(X^2-2X\,\mathbb EX+(\mathbb EX)^2\big) = \mathbb EX^2 - 2\,\mathbb EX\,\mathbb EX + (\mathbb EX)^2 = \mathbb EX^2-(\mathbb EX)^2. \qedbg
\]
\end{proofbox}

\section{Свойства на дисперсията}

\begin{thmbox}[]{Теорема}
Нека $X,Y$ са случайни величини с~$DX<\infty$, $DY<\infty$.
\begin{enumerate}[label=\arabic*)]
  \item \textbf{Неотрицателност.} $DX\ge0$; частен случай: $DX=0 \iff X$ е константа с~вероятност~1.
  \item $D(cX) = c^2 DX$ за константа $c$.
  \item $D(aX+b) = a^2 DX$ за всички $a,b\in\mathbb R$.
  \item Ако $X\perp\!\!\!\perp Y$, то $D(X+Y)=DX+DY$. В общия случай: $D\big(\sum_{i=1}^n X_i\big) = \sum_{i=1}^n DX_i$ (за взаимно независими $X_i$).
\end{enumerate}
\end{thmbox}

\begin{proofbox}
\textbf{1)} $DX=\sum_j(x_j-\mathbb EX)^2 p_j$ - всяко събираемо е
неотрицателно ($(x_j-\mathbb EX)^2\ge0$, $p_j\ge0$), следователно
$DX\ge0$. Ако $X\equiv c$ е константа, то $DC = \mathbb E(c-\mathbb
Ec)^2 = \mathbb E(c-c)^2=0$.

\textbf{2)} $D(cX) = \mathbb E(cX)^2 - (\mathbb EcX)^2 = c^2\mathbb
EX^2 - c^2(\mathbb EX)^2 = c^2\big(\mathbb EX^2-(\mathbb EX)^2\big) =
c^2 DX$.

\textbf{3)} Нека $Y=aX+b$, $\mu=\mathbb EX$, тогава $\mathbb EY=a\mu+b$
и~\[
  DY = \mathbb E\big(aX+b-(a\mu+b)\big)^2 = \mathbb E\big(aX-a\mu\big)^2 = \mathbb E\big(a^2(X-\mu)^2\big) = a^2\,\mathbb E(X-\mu)^2 = a^2 DX.
\]

\textbf{4)}
\begin{align*}
  D(X+Y) &= \mathbb E(X+Y)^2 - \big(\mathbb E(X+Y)\big)^2 = \mathbb E(X^2+2XY+Y^2) - (\mathbb EX+\mathbb EY)^2\\
  &= \mathbb EX^2+\mathbb E(2XY)+\mathbb EY^2 - (\mathbb EX)^2-2\,\mathbb EX\,\mathbb EY-(\mathbb EY)^2\\
  &= DX+DY+2\big(\mathbb E(XY)-\mathbb EX\,\mathbb EY\big).
\end{align*}
Ако $X\perp\!\!\!\perp Y$, по свойство~7 от предната лема $\mathbb
E(XY)=\mathbb EX\,\mathbb EY$, следователно последната скоба е нула
и~$D(X+Y)=DX+DY$. \qedbg
\end{proofbox}

\begin{thmbox}[~(на Марков)]{Неравенство}
Нека $X$ е дискретна случайна величина, $X\ge0$. Тогава за всяко
$\alpha>0$:
\[
  \mathbb P(X>\alpha) \le \frac{\mathbb EX}{\alpha}.
\]
\end{thmbox}

\begin{proofbox}
\[
  \mathbb EX = \sum_{x\in\mathcal X} x\,p(x) \ge \sum_{x\in\mathcal X:\,x>\alpha} x\,p(x) \ge \sum_{x\in\mathcal X:\,x>\alpha} \alpha\,p(x) = \alpha\,\mathbb P(X>\alpha). \qedbg
\]
\end{proofbox}

\begin{thmbox}[~(на Чебишов)]{Неравенство}
Нека $X$ е дискретна случайна величина с~$\mathbb EX<\infty$,
$DX<\infty$. Тогава за всяко $\alpha>0$:
\[
  \mathbb P\big(|X-\mathbb EX|>\alpha\big) \le \frac{DX}{\alpha^2}.
\]
\end{thmbox}

\begin{proofbox}
От неравенството на Марков, приложено към $(X-\mathbb EX)^2\ge0$ и~прага $\alpha^2$:
\[
  \mathbb P\big(|X-\mathbb EX|>\alpha\big) = \mathbb P\big((X-\mathbb EX)^2>\alpha^2\big) \le \frac{\mathbb E\big[(X-\mathbb EX)^2\big]}{\alpha^2} = \frac{DX}{\alpha^2}. \qedbg
\]
\end{proofbox}

\section{Пораждащи функции}

\begin{rembox}[Идея]
Пораждащата (моментна) функция прибира на „торбичка“ или цифров
пакет с~данни цялата вероятностна информация за $X$, запазена и~компресирана в~една удобна за работа математическа форма.
Разгръщайки я - чрез диференциране и~други операции - можем да
извлечем конкретни характеристики на разпределението, като моменти
или вероятности.
\end{rembox}

\begin{defbox}[~(Пораждаща функция)]
Нека $X:\Omega\to\mathbb N$ е случайна величина във $V$. Функцията
\[
  g_X(s) := \mathbb E s^X = \sum_{k=0}^{\infty} s^k\,\mathbb P(X=k), \qquad |s|\le 1,
\]
се нарича \textbf{пораждаща функция} на $X$ (условието $|s|\le1$
гарантира, че редът е сходящ).
\end{defbox}

\begin{thmbox}[]{Теорема}
\begin{itemize}[nosep]
  \item За целочислена случайна величина $X$ с~$\mathbb EX<\infty$: $g_X'(1) = \mathbb EX$.
  \item Ако $\mathbb EX^2<\infty$, то $DX = g_X''(1) + g_X'(1) - (g_X'(1))^2$.
  \item Във всички случаи: $\mathbb P(X=n) = \dfrac{g_X^{(n)}(0)}{n!}$, $\forall n\ge0$.
\end{itemize}
\end{thmbox}

\begin{proofbox}
Диференцираме почленно (термин по термин) реда, определящ $g_X$:
\[
  g_X'(s) = \frac{d}{ds}\,\mathbb Es^X = \mathbb E\frac{d}{ds}s^X = \mathbb EXs^{X-1}.
\]
При $s=1$: $g_X'(1)=\mathbb EX\cdot1=\mathbb EX$.

Диференцираме отново:
\[
  g_X''(s) = \frac{d}{ds}\,\mathbb EXs^{X-1} = \mathbb EX(X-1)s^{X-2},
\]
и~при $s=1$: $g_X''(1)=\mathbb EX(X-1)=\mathbb EX^2-\mathbb EX =
\mathbb EX^2-g_X'(1)$, откъдето $\mathbb EX^2 = g_X''(1)+g_X'(1)$.
Заместваме в~$DX=\mathbb EX^2-(\mathbb EX)^2$ и~получаваме
$DX=g_X''(1)+g_X'(1)-(g_X'(1))^2$.

За последното твърдение, разписваме реда:
\[
  g_X(s) = \sum_{k=0}^{\infty} s^k\,\mathbb P(X=k),
\]
и~вземаме неговата $n$-та производна:
\[
  g_X^{(n)}(s) = \Big(\sum_{k=0}^{\infty}s^k\,\mathbb P(X=k)\Big)^{(n)} = n!\,\mathbb P(X=n) + (n{+}1)n\cdots2\cdot s\,\mathbb P(X=n{+}1) + \dots
\]
(защото членовете с~$k<n$ изчезват при $n$-кратно диференциране).
При $s=0$ всички членове с~положителна степен на $s$ се анулират и~остава $g_X^{(n)}(0) = n!\,\mathbb P(X=n)$, откъдето
$\mathbb P(X=n) = \dfrac{g_X^{(n)}(0)}{n!}$. \qedbg
\end{proofbox}

\begin{thmbox}[]{Лема}
Нека $(X_1,\dots,X_n)$ са независими в~съвкупност целочислени
случайни величини. Тогава
\[
  g_{X_1+\dots+X_n}(s) = \prod_{j=1}^{n} g_{X_j}(s).
\]
Ако в~допълнение случайните величини са еднакво разпределени
($X_1\stackrel{d}{=}X_j$ за всяко $j$, т.е.\ $\mathbb P(X_1=k)=
\mathbb P(X_j=k)$ за всяко $k$), то
\[
  g_{X_1+\dots+X_n}(s) = g_{X_1}^{\,n}(s).
\]
\end{thmbox}

\begin{proofbox}
За $n=2$:
\begin{align*}
  g_{X_1+X_2}(s) &= \mathbb Es^{X_1+X_2} = \sum_{m,n\ge0} s^{m+n}\,\mathbb P(X_1=m,X_2=n)\\
  &\overset{X_1\perp\!\!\!\perp X_2}{=} \sum_{m,n\ge0} s^m s^n\,\mathbb P(X_1=m)\,\mathbb P(X_2=n)\\
  &= \sum_{m=0}^{\infty} s^m\,\mathbb P(X_1=m) \sum_{n=0}^{\infty} s^n\,\mathbb P(X_2=n) = g_{X_1}(s)\,g_{X_2}(s).
\end{align*}
Общият случай за $n$ величини следва по индукция по същата схема.

Ако допълнително $X_1\stackrel{d}{=}X_j$, то $\mathbb P(X_1=n)=
\mathbb P(X_j=n)$ за всяко $n\ge0$, следователно $g_{X_j}(s)=
\sum_n s^n\mathbb P(X_j=n) = \sum_n s^n\mathbb P(X_1=n)=g_{X_1}(s)$
за всяко $j$, откъдето $\prod_{j=1}^n g_{X_j}(s) = g_{X_1}^{\,n}(s)$.
\qedbg
\end{proofbox}

\begin{thmbox}[]{Теорема}
Нека $g_X(s)=g_Y(s)$ за $|s|<1$. Тогава $X\stackrel{d}{=}Y$.
\end{thmbox}

\begin{proofbox}
От предната теорема, $\mathbb P(X=n)=\dfrac{g_X^{(n)}(0)}{n!}$ и~$\mathbb P(Y=n)=\dfrac{g_Y^{(n)}(0)}{n!}$. Понеже самите функции са
равни ($g_X(s)=g_Y(s)$ като функции), то и~производните им от всеки
ред са равни ($g_X^{(n)}(0)=g_Y^{(n)}(0)$ за всяко $n$), следователно
$\mathbb P(X=n)=\mathbb P(Y=n)$ за всяко $n\ge0$, т.е.\ $X$ и~$Y$ са
еднакво разпределени ($X\stackrel{d}{=}Y$). \qedbg
\end{proofbox}
